\slajdid{08-s043}
\begin{frame}[label=08-s043]
\gusttekst
Smatrajmo da je pre trenutka $t_0$ stanje napona na ulaz $U_0$ „beskonačno“ trajalo. Intuitivno, ako bi napona na kondezatoru bio manji od napona $U_0$ postojala bi struja kroz kondenzator i on bi se punio i nekom trenutku bi možda i postao veći od napona $U_0$. A ako bi napon na kondezatoru bio veći od napona $U_0$ postojala bi struja kroz kondenzator suprotnog smera i on bi se praznio, pa bi u nekom trenutku možda bio i manji od napona $U_0$. Međutim bez obzira na smer promene u nekom trenutku napon na kondenzatoru može postati jednak $U_0$. U tom trenutku važi
\[\lim_{u_C(t)\to U_0}i_R(t)=\lim_{u_C(t)\to U_0}i_C(t)=\lim_{u_C(t)\to U_0}\frac{u_C(t)-u_o(t)}R=\lim_{u_C(t)\to U_0}\frac{u_C(t)-U_0}R=0\]
odnosno struja kroz kondenzator postaje jednaka nuli, a tada važi
\[i_C(t)=0\ =>\ \frac{du_C(t)}{dt}=0\ =>\ u_C(t)=const\]
Znači ako je vreme pre trenutka $t_0$ dovoljno dugo trajalo (beskonačno) kolo će ući u stacionarno stanje u kome nema više promena i karakterisano je
\[i_C(t)=0,\qquad u_C(t)=const\]
\end{frame}
